% =========================================================
% Mediants, Continued Fractions, and Rational Approximation
% =========================================================

\subsection{Mediants, Continued Fractions, and Rational Approximation}

% ---------------------------------------------------------
% Teacher exposition
% ---------------------------------------------------------

\begin{remark}[Teacher exposition]
\label{rem:mediant-cf-teacher}

The mediant construction generates rational numbers by repeatedly
refining intervals.  Starting with two fractions

\[
\frac{a}{b} < \frac{c}{d},
\]

the mediant

\[
\frac{a+c}{b+d}
\]

lies strictly between them.  Iterating this process produces infinitely
many rational numbers.

This idea leads naturally to the \emph{Stern--Brocot tree}, a binary tree
that contains every positive rational number exactly once in reduced
form.

A closely related representation of real numbers is given by
\emph{continued fractions}.  Continued fractions generate rational
approximations that are closely connected to the mediant construction.

Together these ideas explain why rational numbers approximate real
numbers extremely well.

\end{remark}

% ---------------------------------------------------------
% Continued fractions
% ---------------------------------------------------------

\begin{definition}[Continued fraction]
\label{def:continued-fraction}

A \emph{continued fraction} is an expression of the form

\[
x = a_0 +
\cfrac{1}{a_1 +
\cfrac{1}{a_2 +
\cfrac{1}{a_3 + \cdots}}}
\]

where $a_0 \in \mathbb{Z}$ and $a_1,a_2,a_3,\dots \in \mathbb{N}$.

\end{definition}

\begin{remark}

Truncating a continued fraction produces rational numbers called
\emph{convergents}.  These convergents provide very accurate rational
approximations to the original number.

\end{remark}

% ---------------------------------------------------------
% Example
% ---------------------------------------------------------

\begin{remark}[Example: continued fraction for $\sqrt{2}$]

The continued fraction expansion of $\sqrt{2}$ is

\[
\sqrt{2} =
1 +
\cfrac{1}{2 +
\cfrac{1}{2 +
\cfrac{1}{2 + \cdots}}}.
\]

The convergents are

\[
1,\quad
\frac{3}{2},\quad
\frac{7}{5},\quad
\frac{17}{12},\quad
\frac{41}{29},\dots
\]

Each convergent is a rational number that approximates $\sqrt{2}$
extremely well.

\end{remark}

% ---------------------------------------------------------
% Connection to mediants
% ---------------------------------------------------------

\begin{remark}[Connection with mediants]

The Stern--Brocot tree generates rational numbers by repeatedly
inserting mediants between neighboring fractions.

Remarkably, the sequence of fractions obtained when searching for a
real number within the Stern--Brocot tree coincides with the sequence
of convergents of the continued fraction expansion of that number.

Thus mediants and continued fractions describe the same approximation
process from two different perspectives.

\end{remark}

% ---------------------------------------------------------
% Diophantine approximation
% ---------------------------------------------------------

\begin{remark}[Diophantine approximation]

Continued fractions provide the best possible rational approximations
to real numbers.

If $\frac{p}{q}$ is a convergent of the continued fraction expansion of
$x$, then

\[
\left| x - \frac{p}{q} \right| < \frac{1}{q^2}.
\]

This estimate is closely related to Dirichlet's Approximation Theorem
and explains why certain rational numbers provide exceptionally good
approximations to real numbers.

\end{remark}

% ---------------------------------------------------------
% Conceptual summary
% ---------------------------------------------------------

\begin{remark}[Conceptual summary]

Several important ideas are unified by these constructions:

\begin{itemize}

\item The mediant operation explains the density of the rational numbers.

\item The Stern--Brocot tree generates every positive rational number.

\item Continued fractions produce optimal rational approximations to
real numbers.

\item These approximations form sequences that converge to real numbers,
revealing the gaps present in the rational number system.

\end{itemize}

Together these ideas connect number theory with the foundations of
analysis.

\end{remark}